\bookchapter{Bab 30: Screen Mirroring, Presentasi, dan Koneksi Mobil}

\chapterintro{Layar ponsel sebesar apa pun cuma 7 inci, tapi TV Anda 55 inci, proyektor Anda 100 inci. Melempar tampilan ponsel ke layar besar, pengalaman berbeda dimensi. Bab ini menjelaskan screen casting, presentasi, dan car connectivity tiga skenario frekuensi tinggi, membantu memperluas kemampuan ponsel ke layar yang lebih besar.}

\booksection{Tiga Kubu Screen Cast Nirkabel — Jangan Tertukar}

Screen cast nirkabel tampak semua “tampilan ponsel ke TV”, tapi teknologi dasarnya ada tiga, tidak saling kompatibel:

\begin{center}
\begin{tabularx}{\linewidth}{|X|X|X|}
\hline
\textbf{Teknologi} & \textbf{Siapa pakai} & \textbf{Karakteristik} \tabularnewline
\hline
AirPlay & iPhone/iPad/Mac & Ekosistem Apple, pengalaman paling lancar. Butuh Apple TV atau TV yang mendukung AirPlay \tabularnewline
Miracast & Ponsel Android & Sambungan langsung Wi-Fi Direct, tidak butuh router. "Screen cast" ponsel Android default pakai ini \tabularnewline
DLNA & Semua platform & Hanya push file media (video/musik/foto), tidak mirror layar. Cocok untuk push video ke TV untuk diputar \tabularnewline
\hline
\end{tabularx}
\end{center}

\booksubsection{Cara operasi masing-masing platform:}

\begin{itemize}
  \item \textbf{: Pengaturan Nama Anda iCloud iCloud Cloud Backup Aktifkan. Ponsel terhubung daya, terhubung Wi-Fi, saat layar terkunci akan otomatis dicadangkan. 5GB gratis kemungkinan besar tidak cukup, disarankan upgrade ke 50GB dengan 6 yuan.}: Control Center > Screen Mirroring > Pilih TV/Apple TV Anda. Jika TV tidak mendukung AirPlay, bisa beli Apple TV box atau TV stick yang mendukung AirPlay.
  \item \textbf{Android}: Tarik bawah notification bar > Cari tombol "Wireless Cast" / "Smart Cast" / "Multi-Screen Interaction" > Pilih TV. Ponsel dan TV harus terhubung Wi-Fi yang sama.
  \item \textbf{Solusi universal}: Jika TV terlalu tua tidak mendukung screen cast nirkabel apa pun, beli Chromecast atau Mi Box/Huawei Box colok ke HDMI TV, bisa menerima cast.
\end{itemize}

\booksection{Screen Cast Berkabel — Lebih Stabil, Zero Latency}

Screen cast nirkabel praktis tapi ada delay, main game dan presentasi pengalaman buruk. Screen cast berkabel bisa selesaikan masalah ini:

\begin{itemize}
  \item \textbf{: Pengaturan Nama Anda iCloud iCloud Cloud Backup Aktifkan. Ponsel terhubung daya, terhubung Wi-Fi, saat layar terkunci akan otomatis dicadangkan. 5GB gratis kemungkinan besar tidak cukup, disarankan upgrade ke 50GB dengan 6 yuan.}: Butuh kabel konverter Lightning/USB-C to HDMI (sekitar 100-200 yuan). Setelah dicolok tampilan ponsel langsung output ke HDMI TV/proyektor.
  \item \textbf{Android}: Sebagian ponsel mendukung USB-C to HDMI langsung (ponsel harus mendukung DP Alt Mode, sebagian besar flagship mendukung). Beli kabel USB-C to HDMI saja.
  \item \textbf{Skenario ruang rapat}: Jika sering perlu presentasi dengan ponsel, disarankan selalu bawa kabel HDMI di tas. TV/proyektor ruang rapat hampir semua punya port HDMI, colok langsung pakai, jauh lebih andal daripada repot dengan screen cast nirkabel.
\end{itemize}

\booksection{Ponsel Jadi Komputer — Mode Desktop}

Sebagian ponsel Android mendukung “mode desktop” — setelah sambung ke monitor, ponsel akan output antarmuka mirip desktop komputer, bisa operasi dengan mouse keyboard:

\begin{itemize}
  \item \textbf{Samsung DeX}: Fitur andalan flagship Samsung. Setelah sambung ke monitor muncul antarmuka desktop, bisa operasi multi-window, pengalaman paling mendekati komputer sungguhan. Dengan keyboard mouse Bluetooth, bisnis keluar kota tidak perlu bawa laptop.
  \item \textbf{Mode Desktop Huawei}: Seri Huawei Mate/Pura setelah sambung ke monitor juga punya fitur serupa. Mendukung multi-window, file management, App office.
  \item \textbf{Motorola Ready For}: Flagship Moto juga mendukung, fitur mirip DeX.
  \item \textbf{: Produksi profesional kelas atas. Seri iPhone 15 Pro mendukung shooting video Apple ProRes (hingga 4K 60fps), tetapi format ProRes RAW saat ini perlu digunakan bersama perangkat perekam eksternal, ponsel belum mendukung shooting langsung.}: Light office (balas email, edit dokumen), presentasi PPT, lihat video di layar besar. Tidak cocok untuk operasi Excel berat atau desain profesional.
\end{itemize}

\booksection{Car Connectivity — Cara Benar Pakai Ponsel Saat Berkendara}

Mengoperasikan ponsel saat berkendara itu melanggar hukum, tapi navigasi, dengar musik, telepon adalah kebutuhan. Sistem car connectivity membuat Anda menggunakan fungsi inti ponsel di layar head unit, ponsel bisa diletakkan tidak disentuh:

\begin{itemize}
  \item \textbf{Apple CarPlay}: Pilihan utama pengguna iPhone. Melalui USB atau nirkabel sambung ke head unit, layar head unit muncul antarmuka sederhana iPhone — navigasi (Amap/Baidu), musik, telepon, pesan, podcast. Kontrol suara Siri sepenuhnya tanpa tangan.
  \item \textbf{Android Auto}: Solusi corresponden pengguna Android. Fitur mirip CarPlay, mendukung Google Maps, Spotify dll. Di dalam negeri karena pembatasan layanan Google, sebagian fitur tidak tersedia.
  \item \textbf{Huawei HiCar}: Solusi buatan Huawei, tingkat dukungan model dalam negeri tinggi. Selain navigasi musik, juga bisa pakai panggilan video Huawei, fusi gambar parkir dan fitur khas lainnya.
  \item \textbf{Baidu CarLife}: Solusi cross-platform, iOS dan Android bisa pakai. Fitur lebih sedikit tapi kompatibilitas luas, banyak model lama hanya mendukung CarLife.
\end{itemize}

\begin{tcbwarn}{Peringatan Keamanan}
Apa pun solusi car connectivity, saat berkendara harus semaksimal mungkin pakai kontrol suara. Navigasi setel tujuan sebelumnya, musik pilih playlist sebelumnya, jangan operasi layar head unit saat berkendara.

\end{tcbwarn}

\booksection{Praktik Screen Cast Ruang Rapat — Tidak Lagi Panik}

Rapat adalah skenario screen cast paling sering, juga skenario paling mudah bermasalah. Berikut beberapa pengalaman praktik:

\begin{itemize}
  \item \textbf{Datang lebih awal untuk tes}: Jangan pernah tunggu rapat mulai baru sambung screen cast. Datang 10 menit lebih awal, tes cara koneksi. Probabilitas screen cast Wi-Fi gagal jauh lebih tinggi dari yang Anda bayangkan.
  \item \textbf{Siapkan solusi berkabel sebagai Plan B}: Taruh kabel HDMI di tas. Saat screen cast nirkabel tidak bisa sambung, kabel adalah penyelamat.
  \item \textbf{Matikan preview notifikasi}: Saat screen cast pesan WeChat Anda akan tampil di layar besar ruang rapat. Sebelum cast matikan "preview konten notifikasi lock screen" di pengaturan, atau aktifkan "mode presentasi" (sebagian ponsel mendukung, saat cast otomatis sembunyikan notifikasi).
  \item \textbf{Pakai "dokumen presentasi" bukan "operasi ponsel"}: Jika ingin menampilkan PPT, langsung simpan PPT di ponsel buka dengan WPS full screen, bukan cast seluruh desktop ponsel. Lebih profesional dan lancar.
\end{itemize}

\begin{tcbpractice}{Praktik: Selesaikan satu kali screen cast ponsel ke TV}
\begin{enumerate}
  \item Di TV temukan fungsi "Screen Mirroring" atau "AirPlay" dan aktifkan.
  \item Buka satu foto atau video di ponsel.
  \item Coba screen cast nirkabel (iPhone pakai Screen Mirroring, Android pakai tombol wireless cast).
  \item Jika gagal, cek apakah TV mendukung protokol yang sesuai, atau coba DLNA push (di App galeri/video tekan tombol "Cast").
  \item Setelah cast berhasil, coba lakukan operasi lain di ponsel, lihat apakah layar besar tampil sinkron.
\end{enumerate}

\end{tcbpractice}

\begin{tcbsummary}{Ringkasan Bab}
\begin{itemize}
  \item Tiga kubu screen cast nirkabel: AirPlay (Apple), Miracast (Android), DLNA (cross-platform push media)
  \item Screen cast berkabel lebih stabil zero latency, skenario ruang rapat wajib kabel HDMI
  \item Samsung DeX / mode desktop Huawei mengubah ponsel jadi komputer light office
  \item Car connectivity: CarPlay / Android Auto / HiCar / CarLife empat solusi besar
  \item Aturan wajib screen cast ruang rapat: tes lebih awal, siapkan Plan B, matikan preview notifikasi
\end{itemize}

\end{tcbsummary}
